
\subsubsection{4.1 Research Questions}

\textbf{RQ1 (Clustering Structure):} Do time series trajectories partition into distinct clusters with silhouette score > 0.25, optimal k in [3,8], and bootstrap stability > 0.65?

\textbf{RQ2 (Phase Detection):} Do trajectories exhibit statistically significant changepoints with detection rate > 50\%?

\textbf{RQ3 (Shape Differentiation):} Do identified clusters represent different trajectory shapes as measured by shape descriptor variance ratios > 2.0?

\textbf{RQ4 (Archetype Recovery):} Do empirical clusters map to proposed behavioral archetypes with $\geq 3$ of 5 archetypes recovered at >70\% alignment?

\subsubsection{4.2 Dataset}

The methodology was validated on 500 time series, each with 300 temporal observations. Data were sourced from the helenqu/astro-time-series dataset on HuggingFace Hub, loaded via the HuggingFace datasets library in streaming mode.

\textbf{Data source clarification:} The original experimental design specified HuggingFace dataset download statistics. However, the HuggingFace Hub API does not expose historical monthly download time series---only current total download counts. To validate the methodology with real time series data from the HuggingFace ecosystem, astronomical lightcurve measurements were used as proxy data.

\subsubsection{4.3 Baselines}

Three baseline approaches were evaluated:

\begin{enumerate}
\item \textbf{Random Assignment:} Random cluster labels (null baseline)
\item \textbf{K-Means on Summary Statistics:} Clustering on aggregate features (mean, standard deviation, trend slope)
\item \textbf{Single-Level DTW:} DTW clustering without PELT phase detection
\end{enumerate}

\subsubsection{4.4 Implementation}

\begin{itemize}
\item Clustering: tslearn v0.6.0 TimeSeriesKMeans with DTW metric
\item Changepoint detection: ruptures v1.1.7 PELT algorithm
\item Evaluation: scikit-learn v1.0 for silhouette scores
\end{itemize}

\textbf{Hyperparameters:}

\begin{table}[htbp]
\centering
\resizebox{\columnwidth}{!}{%
\begin{tabular}{ll}
\toprule
Parameter & Value \\
\midrule
k range & [3, 8] \\
DTW metric & Full DTW \\
PELT penalty & BIC = 2log(n) \\
Bootstrap iterations & 100 \\
Bootstrap sample ratio & 80\% \\
\bottomrule
\end{tabular}
}
\end{table}
